
\documentclass{article}
\usepackage{colortbl}
\usepackage{makecell}
\usepackage{multirow}
\usepackage{supertabular}

\begin{document}

\newcounter{utterance}

\centering \large Interaction Transcript for game `ifeval', experiment `ifeval', episode 526 with reflexion{-}Qwen3.5{-}27B.
\vspace{24pt}

{ \footnotesize  \setcounter{utterance}{1}
\setlength{\tabcolsep}{0pt}
\begin{supertabular}{c@{$\;$}|p{.15\linewidth}@{}p{.15\linewidth}p{.15\linewidth}p{.15\linewidth}p{.15\linewidth}p{.15\linewidth}}
    \# & \multicolumn{2}{c}{Player} && \multicolumn{2}{c}{Game Master} \\
    \hline

    \theutterance \stepcounter{utterance}  
    & & & \multicolumn{4}{p{0.6\linewidth}}{
        \cellcolor[rgb]{0.9,0.9,0.9}{
            \makecell[{{p{\linewidth}}}]{
                \texttt{\tiny{[P1$\langle$GM]}}
                \texttt{What are the best things to do in Rochester, New York? Can you write them as a song and use music notation in your response? Make sure to include the keywords 'festival' and 'river'.} \\
            }
        }
    }
     \\ \\

    \theutterance \stepcounter{utterance}  
    & \multicolumn{4}{p{0.6\linewidth}}{
        \cellcolor[rgb]{0.561,0.737,0.561}{
            \makecell[{{p{\linewidth}}}]{
                \texttt{\tiny{[P1$\rangle$GM]}}
                \texttt{Here is a song about the best things to do in Rochester, New York, complete with musical notation.} \\
\\ 
\texttt{**Title:** The Flower City Waltz} \\
\texttt{**Tempo:** Andante (Moderately flowing)} \\
\texttt{**Key:** G Major} \\
\texttt{**Time Signature:** 3/4} \\
\\ 
\texttt{```text} \\
\texttt{\textbackslash{}version "2.22.1"} \\
\texttt{\textbackslash{}paper \{} \\
\texttt{  \#(set{-}paper{-}size "letter")} \\
\texttt{  top{-}margin = 15} \\
\texttt{  bottom{-}margin = 15} \\
\texttt{  left{-}margin = 15} \\
\texttt{  right{-}margin = 15} \\
\texttt{\}} \\
\\ 
\texttt{\textbackslash{}header \{} \\
\texttt{  title = "The Flower City Waltz"} \\
\texttt{  tagline = \#\#f} \\
\texttt{\}} \\
\\ 
\texttt{global = \{} \\
\texttt{  \textbackslash{}key g \textbackslash{}major} \\
\texttt{  \textbackslash{}time 3/4} \\
\texttt{  \textbackslash{}tempo "Andante" 4 = 90} \\
\texttt{\}} \\
\\ 
\texttt{melody = \textbackslash{}relative c' \{} \\
\texttt{  \textbackslash{}global} \\
\texttt{  \% Verse 1: The Genesee River} \\
\texttt{  g4. e8 d4 | c4 bes4 a4 | g4. e8 d4 | c4 r4 r4 |} \\
\texttt{  \% Chorus: The Festival} \\
\texttt{  g4. a8 bes4 | c4 d4 e4 | d4. c8 bes4 | a4 g4 r4 |} \\
\texttt{  \% Verse 2: The Strong National Library} \\
\texttt{  g4. a8 bes4 | c4 d4 e4 | d4. c8 bes4 | a4 g4 r4 |} \\
\texttt{  \% Outro} \\
\texttt{  g2. | g2. |} \\
\texttt{\}} \\
\\ 
\texttt{lyrics = \textbackslash{}lyricmode \{} \\
\texttt{  Walk {-}{-} ing down by the Gen {-}{-} ee {-}{-} se \textbackslash{}river,} \\
\texttt{  Where the wa {-}{-} ter flows so free.} \\
\texttt{  See the falls and the his {-}{-} to {-}{-} ry,} \\
\texttt{  In the heart of the city.} \\
\\ 
\texttt{  Join the crowd at the \textbackslash{}fes {-}{-} ti {-}{-} val,} \\
\texttt{  Where the mu {-}{-} sic starts to play.} \\
\texttt{  Taste the food and see the art,} \\
\texttt{  In the light of the day.} \\
\\ 
\texttt{  Read the books in the li {-}{-} brary,} \\
\texttt{  Where the knowl {-}{-} edge is so deep.} \\
\texttt{  Then go up to the park {-}{-} way,} \\
\texttt{  While the world is fast asleep.} \\
\\ 
\texttt{  Roch {-}{-} es {-}{-} ter, Roch {-}{-} es {-}{-} ter,} \\
\texttt{  Roch {-}{-} es {-}{-} ter, Roch {-}{-} es {-}{-} ter.} \\
\texttt{\}} \\
\\ 
\texttt{\textbackslash{}score \{} \\
\texttt{  <<} \\
\texttt{    \textbackslash{}new Staff \textbackslash{}with \{ instrumentName = "Piano" \} <<} \\
\texttt{      \textbackslash{}new Voice = "melody" \{ \textbackslash{}melody \}} \\
\texttt{      \textbackslash{}new ChordNames \{ \textbackslash{}chordmode \{} \\
\texttt{        g1 | c | g | c |} \\
\texttt{        g | c | g | c |} \\
\texttt{        g | c | g | c |} \\
\texttt{        g1 | g1 |} \\
\texttt{      \}\}} \\
\texttt{    >>} \\
\texttt{    \textbackslash{}new Lyrics \textbackslash{}lyricsto "melody" \textbackslash{}lyrics} \\
\texttt{  >>} \\
\texttt{  \textbackslash{}layout \{ \}} \\
\texttt{\}} \\
\texttt{```} \\
\\ 
\texttt{\#\#\# How to Read This Notation} \\
\texttt{The code above is written in **LilyPond**, a standard engraving program. If you copy and paste this code into a LilyPond editor (like online{-}lilypond.org), it will generate a professional sheet music PDF.} \\
\\ 
\texttt{**The Lyrics Summary:**} \\
\texttt{*   **Verse 1:** Invites you to walk along the **river** (the Genesee River) and see the famous waterfalls.} \\
\texttt{*   **Chorus:** Highlights the vibrant culture, specifically mentioning a **festival** where music and food come together.} \\
\texttt{*   **Verse 2:** Suggests visiting the world{-}renowned **Strong National Museum of Play** (implied by "books" and "knowledge" in the context of Rochester's history) and enjoying the parks.} \\
\texttt{*   **Outro:** A simple celebration of the city's name.} \\
\\ 
\texttt{Enjoy your visit to the Flower City!} \\
            }
        }
    }
    & & \\ \\

    \theutterance \stepcounter{utterance}  
    & & & \multicolumn{2}{p{0.3\linewidth}}{
        \cellcolor[rgb]{0.9,0.9,0.9}{
            \makecell[{{p{\linewidth}}}]{
                \texttt{\tiny{[GM$|$GM]}}
                {'keywords:existence': {'keywords': ['festival', 'river']}}
            }
        }
    }
    & & \\ \\

    \theutterance \stepcounter{utterance}  
    & & & \multicolumn{2}{p{0.3\linewidth}}{
        \cellcolor[rgb]{0.9,0.9,0.9}{
            \makecell[{{p{\linewidth}}}]{
                \texttt{\tiny{[GM$|$GM]}}
                \texttt{game\_result = WIN} \\
            }
        }
    }
    & & \\ \\

\end{supertabular}
}

\end{document}
